\chapter{A Rule-Based Route to Clause Normal Form}
\label{sec:p12-small-clause-standard}

Clausification converts a first-order problem into the clause language
used by many deduction procedures. This chapter gives the standard
seven-stage procedure in Nonnengart and Weidenbach's presentation,
including its treatment of polarity, variables and clauses
\cite[secs.~1--3, pp.~337--347]{NonnengartWeidenbach2001SmallClauses}.
The later formula-renaming and optimized-Skolemization methods have
their own constructions. Clause count, literal count, term size and
the time taken by a particular prover are different measures.
The preservation results establish logical guarantees, not a
universal ordering of prover runtimes.

\section{Input semantics and the objects being transformed}

\begin{hbdefinition}{def:p12-small-language}{Language, interpretations and clauses}
Use single-sorted first-order logic over a signature
\(\Sigma=(\mathcal F,\mathcal R)\) of function and predicate symbols
with fixed finite arities and an infinite supply of variable names \(X\).
Terms are variables, constants and applications \(f(t_1,\ldots,t_n)\).
Atoms are \(R(t_1,\ldots,t_n)\); equality, if included, is an atom
interpreted as identity. Formula constructors are
\(\top,\bot,\neg,\land,\lor,\to,\leftrightarrow,\forall,\exists\).
The source's \(\supset,\equiv\) denote our object-language
\(\to,\leftrightarrow\).

A structure \(\mathfrak A\) has a nonempty domain \(D\), an
interpretation of every signature symbol, and valuations \(\nu:X\to D\).
The updated valuation \(\nu[x/d]\) sends \(x\) to \(d\) and
agrees with \(\nu\) elsewhere. Write \(t^{\mathfrak A,\nu}\)
for term evaluation: variables take their valuation, constants their
specified interpretation, and
\[
 f(t_1,\ldots,t_n)^{\mathfrak A,\nu}
 =f^{\mathfrak A}(t_1^{\mathfrak A,\nu},\ldots,t_n^{\mathfrak A,\nu}).
\]
An atom \(R(\bar t)\) is true exactly when the tuple of term values
belongs to \(R^{\mathfrak A}\). The constants \(\top,\bot\)
are respectively true and false; negation reverses truth, conjunction
requires both operands, disjunction requires at least one, implication
is false exactly when its antecedent is true and its consequent false,
and a biconditional requires equal truth values. Universal quantification
requires the body under every \(\nu[x/d]\), existential quantification
under at least one, for \(d\in D\).
Write \(\Phi\equiv_\Sigma\Psi\)
for equality of truth in every \(\Sigma\)-structure and every valuation.
In displayed equivalence chains, \(\equiv\) abbreviates this relation
for the current fixed signature; a Skolem step is identified separately
as an expansion/reduct transformation.
Satisfiability of an open formula means that some structure and
\emph{some valuation} satisfy it.

A literal is an atom or a negated atom. A clause is a finite
\emph{multiset} of literals, read disjunctively with all its variables
universally quantified. A clause set is read conjunctively.
The empty clause is false; the empty clause set is true.
The formula's free variables, bound variable names and their union
are denoted \(\operatorname{free},\operatorname{bound},\operatorname{vars}\).
An atom's free variables are the union of the variables in its terms,
and it has no bound variables. Logical constants have neither.
Negation inherits both sets from its operand.
Free variables are obtained by union at binary connectives and deletion
of the bound name at a quantifier; bound names are accumulated at
quantifiers. These are sets of names, distinct from binder occurrences.
A formula is \emph{rectified} when different binder occurrences have
different names and no bound name is also free in the formula.
\end{hbdefinition}

\begin{hbtheorem}{thm:p12-small-sat-closure}{proved=annex:p12.small.sat.closure}{The open-input convention}
If \(\operatorname{free}(\Phi)=\{x_1,\ldots,x_k\}\), put
\[
                   \Phi^\exists=\exists x_1\cdots\exists x_k\,\Phi .
\]
For every structure \(\mathfrak A\),
\(\mathfrak A\models\Phi^\exists\) exactly when some valuation
satisfies \(\Phi\). For a sentence this closure changes nothing.
If an input instead stipulates universal closure, apply the procedure
to that explicitly specified sentence. The two conventions cannot
be silently interchanged: \(P(x)\land\neg P(y)\) is satisfiable
under the open-input convention but its universal closure is not.
\end{hbtheorem}

\begin{hbdefinition}{def:p12-small-positions}{Positions, polarity and substitution}
A formula position is a finite word over child indices. The root
is \(\epsilon\); negation and quantifiers have child \(1\), and
binary connectives have children \(1,2\). Atoms have no formula
children. Write \(\Phi|_\pi\) for the selected subformula and
\(\Phi[\pi/\Psi]\) for its replacement. Concatenation defines
the prefix ordering; the length of a position is its number of indices.
Formula depth is the maximum position length. Term depth is defined
separately by depth zero for variables/constants and one plus the
maximum child depth at a positive-arity function.

Set \(\operatorname{pol}(\Phi,\epsilon)=+1\).
Conjunction, disjunction and quantifiers preserve the incoming sign.
An implication preserves it in its consequent and reverses it in
its antecedent; negation reverses it. Below a biconditional the
polarity is \(0\), and subsequent operations keep it \(0\).

A term substitution \(\sigma\) is a simultaneous variable-to-term
map with finite support; its domain is the changed variables and
its changed-image set is the set of their images. Application to
terms is homomorphic. Application to formulas in witness substitution
is capture-avoiding. The separate raw variable-renaming operation
below changes binder names as well as variable occurrences.
\end{hbdefinition}

\begin{hblemma}{thm:p12-small-context}{proved=annex:p12.small.context}{Replacement at an occurrence}
Equivalent subformulas can be replaced at any occurrence, including
under binders. More generally, suppose \(A\) entails \(B\) under
every valuation in the structures being considered. Replacing
an occurrence of \(A\) by \(B\) preserves implication from the old
whole formula to the new one at positive polarity, and reverses
that implication at negative polarity. At polarity zero,
equivalence of the inserted formulas is sufficient.
These statements compare the same interpretation of all symbols
and allow bound variables to take every domain value.
\end{hblemma}

\section{Simplification and negation normal form}

\begin{hbmethod}{def:p12-small-simplification}{The complete initial simplification rules}
Apply these rules at subformula occurrences; match
\(\land,\lor,\leftrightarrow\) modulo commutativity.
Matching is not itself an extra rewrite step.
\[
\begin{array}{lll@{\qquad}lll}
A\land A&\longmapsto&A&A\lor A&\longmapsto&A\\
A\land\neg A&\longmapsto&\bot&A\lor\neg A&\longmapsto&\top\\
A\leftrightarrow A&\longmapsto&\top&A\to A&\longmapsto&\top\\
\neg\top&\longmapsto&\bot&\neg\bot&\longmapsto&\top\\
A\land\top&\longmapsto&A&A\lor\top&\longmapsto&\top\\
A\land\bot&\longmapsto&\bot&A\lor\bot&\longmapsto&A\\
A\to\top&\longmapsto&\top&\top\to A&\longmapsto&A\\
A\to\bot&\longmapsto&\neg A&\bot\to A&\longmapsto&\top\\
A\leftrightarrow\top&\longmapsto&A&
A\leftrightarrow\bot&\longmapsto&\neg A .
\end{array}
\]
Also replace \(Qx\,A\) by \(A\), for either quantifier \(Q\),
when \(x\notin\operatorname{free}(A)\).
\end{hbmethod}

\begin{hbtheorem}{thm:p12-small-simplify}{proved=annex:p12.small.simplify}{Termination and the simplification invariant}
Every simplification step preserves equivalence and strictly reduces
the formula-tree node count. Every exhaustive sequence terminates.
Its result is \(\top\), \(\bot\), or a formula containing neither
logical constant, with every remaining quantifier binding a variable
free in its body. Free variables may disappear during simplification.
If the result is a logical constant, return the corresponding
empty clause set or singleton empty-clause set.
\end{hbtheorem}

\begin{hbmethod}{def:p12-small-nnf-rules}{Polarity-dependent NNF rewriting}
At a biconditional occurrence \(A\leftrightarrow B\), use
\[
\begin{array}{ll}
(A\to B)\land(B\to A),&\operatorname{pol}=+1,\\
(A\land B)\lor(\neg A\land\neg B),&\operatorname{pol}=-1.
\end{array}
\]
The remaining rules are
\[
\begin{aligned}
\neg(A\land B)&\longmapsto\neg A\lor\neg B,&
\neg(A\lor B)&\longmapsto\neg A\land\neg B,\\
\neg\forall x\,A&\longmapsto\exists x\,\neg A,&
\neg\exists x\,A&\longmapsto\forall x\,\neg A,\\
A\to B&\longmapsto\neg A\lor B,&
\neg\neg A&\longmapsto A .
\end{aligned}
\]
An outermost remaining biconditional always has nonzero polarity.
No rule need operate on a biconditional below another one.
\end{hbmethod}

\begin{hbtheorem}{thm:p12-small-nnf}{proved=annex:p12.small.nnf}{Every permitted NNF rewrite sequence terminates}
On a constant-free input, these rules terminate in an equivalent NNF:
only \(\land,\lor,\forall,\exists\) connect literals.
They preserve both the free-variable and all-variable sets.
A strictly decreasing interpretation is obtained from atomic weight \(1\):
negation has weight \(2a\); either quantifier has weight \(a+1\);
\(\land,\lor\) have weight \(a+b+1\); implication has weight
\(3a+2b+2\); biconditional has weight \(6a+6b+6\).
The proof permits duplication of complex subformulas.

For atomic \(A,B\), polarity-dependent treatment of
\(\neg(A\leftrightarrow B)\) gives the two clauses
\(A\lor B\) and \(\neg A\lor\neg B\).
Always using the positive expansion instead also generates
\(A\lor\neg A\) and \(B\lor\neg B\) before simplification.
\end{hbtheorem}

\section{Miniscoping and variable rectification}

\begin{hbmethod}{def:p12-small-miniscoping}{The six miniscoping schemes}
For \(Q\in\{\forall,\exists\}\) and
\(\circ\in\{\land,\lor\}\), use
\[
            Qx(A\circ B)\longmapsto(Qx A)\circ B
                \quad\text{if }x\notin\operatorname{free}(B).
\]
These are four rules, also matched with the operands exchanged.
The other two are
\[
\begin{aligned}
\forall x(A\land B)&\longmapsto(\forall x A)\land(\forall x B),\\
\exists x(A\lor B)&\longmapsto(\exists x A)\lor(\exists x B),
\end{aligned}
\]
when \(x\) occurs free in both operands.
No corresponding unrestricted distribution of \(\forall\) over
\(\lor\), or \(\exists\) over \(\land\), is allowed.
\end{hbmethod}

\begin{hbtheorem}{thm:p12-small-miniscope}{proved=annex:p12.small.miniscope}{Miniscoping preserves truth and terminates}
On NNF formulas every displayed step preserves equivalence, and
every rewrite sequence terminates. Let \(b(A)\) count the
\(\land,\lor\) occurrences in \(A\). The nonnegative integer
\[
          \sum_{\text{quantifier occurrences }Qx A}3^{b(A)}
\]
strictly decreases at each step, even when a quantifier is duplicated.
It is zero when the formula is quantifier-free.
The objective is to expose smaller witness-dependency lists.
This does not promise a globally minimal number of Skolem symbols
or a faster proof search.
\end{hbtheorem}

\begin{hbdefinition}{def:p12-small-raw-renaming}{Raw renaming on the occurring variables}
For an injective map
\(\rho:\operatorname{vars}(A)\to X\), define
\(\operatorname{raw}_\rho(A)\) by replacing every variable
occurrence and every binder name \(x\) by \(\rho(x)\).
Use the same homomorphic replacement in terms and preserve
all connectives and quantifier kinds.
For \(y\notin\operatorname{vars}(A)\), the notation
\(\operatorname{raw}_{x\mapsto y}(A)\) changes \(x\) and
fixes the other occurring variables.
Injectivity is required on the variables of the formula being
renamed, not on an identity-extended map of the entire alphabet.
\end{hbdefinition}

\begin{hblemma}{thm:p12-small-renaming}{proved=annex:p12.small.renaming}{The valuation lemma, including both quantifiers}
If \(\nu_1(z)=\nu_2(\rho(z))\) for every
\(z\in\operatorname{free}(A)\), then
\[
       \mathfrak A,\nu_1\models A
       \quad\Longleftrightarrow\quad
       \mathfrak A,\nu_2\models\operatorname{raw}_\rho(A).
\]
For terms whose variables are in the domain of \(\rho\), term values
obey the corresponding equality when the valuations agree modulo
\(\rho\) on those variables. This is the semantic interface used
by the indexed renaming algorithm below.
\end{hblemma}

\begin{hbmethod}{def:p12-small-rectification}{The indexed renaming procedure}
Let \(nq(A)\) count quantifier occurrences. Supply distinct fresh
variables \(y_1,\ldots,y_{nq(A)}\) outside
\(\operatorname{vars}(A)\). Define \(\operatorname{ren}(A,n)\)
recursively, using the appropriate unused portion of this list.
Atoms and logical constants are unchanged; negation recurses
with the same \(n\). At \(B\circ C\), for any binary connective,
return
\[
       \operatorname{ren}(B,n)\circ
       \operatorname{ren}(C,n+nq(B)).
\]
At \(Qx B\), return
\[
          Qy_n\,\operatorname{ren}
                    (\operatorname{raw}_{x\mapsto y_n}(B),n+1).
\]
Start at \(n=1\). The finite list is the portion actually used
of the source's infinite fresh-variable sequence.
\end{hbmethod}

\begin{hbtheorem}{thm:p12-small-rectification}{proved=annex:p12.small.rectification}{Correctness and allocation of fresh names}
The procedure terminates, preserves equivalence and the free-variable
set, and assigns different fresh names to all binder occurrences.
No bound name is free in the resulting formula.
A call starting at \(n\) uses exactly the binder-name block
\(y_n,\ldots,y_{n+nq(A)-1}\), with the empty-block convention
when there is no quantifier.
\end{hbtheorem}

\section{Witness functions and the clause matrix}

\begin{hbmethod}{def:p12-small-skolem-rule}{Outermost local Skolemization}
Start with a closed, rectified NNF. Choose an existential
occurrence \(\exists x B\) of minimum position length.
List \(\operatorname{free}(\exists x B)=\{z_1,\ldots,z_m\}\).
Introduce a function symbol \(f\) of arity \(m\), fresh outside
the current declared signature, and replace that occurrence by
\[
                         B[x:=f(z_1,\ldots,z_m)] .
\]
The replacement concerns the free occurrences of \(x\) in \(B\).
Rectification prevents capture. Arity zero gives a new constant.
Record the new symbol and its ordered dependency list, then repeat.
\end{hbmethod}

\begin{hbtheorem}{thm:p12-small-skolem}{proved=annex:p12.small.skolem}{Local expansion and reduct, on arbitrary nonempty domains}
Every step of the displayed rule removes one existential quantifier.
Assuming Choice, each model of the old closed formula has an
expansion on the same domain satisfying the new formula.
Without Choice, every model of the new formula has a reduct
satisfying the old one. The expansion can make
\(\exists x B\) and \(B[x:=f(\bar z)]\) agree at every valuation.
No finite-domain restriction is imposed.
The outermost schedule leaves only universally bound parameters
above the selected occurrence.

Freshness outside the declared signature gives this full
expansion guarantee. A name merely absent from the input formula
may already have a fixed interpretation in that signature.
Reusing it can require changing that interpretation; this is
a different model relation, even when satisfiability is preserved.
\end{hbtheorem}

\begin{hbmethod}{def:p12-small-clause-matrix}{A schedule that preserves rectification}
First pull all remaining universal quantifiers to a prefix using
\[
          (\forall x A)\circ B\longmapsto\forall x(A\circ B),
       \qquad \circ\in\{\land,\lor\},\quad
                    x\notin\operatorname{free}(B).
\]
Only after the prefix has been formed, distribute in its
quantifier-free matrix:
\[
                 A\lor(B\land C)
                 \longmapsto(A\lor B)\land(A\lor C).
\]
Commutative matching covers either operand order.
Flatten the final conjunction and its disjunctions into clauses.
\end{hbmethod}

\begin{hbtheorem}{thm:p12-small-cnf}{proved=annex:p12.small.cnf}{Prefix formation and distributive termination}
The first phase terminates and produces a universal prefix:
the sum of the numbers of Boolean ancestors of all quantifier
occurrences decreases by one per pull.
Global rectification makes every required side condition available.
The second phase terminates in CNF. Give each literal weight \(2\),
give conjunction weight \(a+b+1\), and disjunction weight \(ab\).
Every distribution strictly decreases this positive integer.
Both phases preserve equivalence in the expanded signature.
Distributing before prefix formation can duplicate binder names
and block further pulling; the proof gives a complete example.
\end{hbtheorem}

\begin{hbdefinition}{def:p12-small-clause-reductions}{Subsumption, tautology and condensation}
For literal multisets, \(C\sigma\subseteq D\) means that every
literal occurs in \(D\) with at least its multiplicity in \(C\sigma\).
Clause \(C\) subsumes \(D\) if some term substitution has this property.
A syntactic tautology contains both an atom and its literal negation.
A condensation \(C\) of \(D\) is obtained from an instance \(D\sigma\)
by deleting duplicate literals, has fewer literal occurrences,
and itself subsumes \(D\).

Delete a tautological clause. Delete a subsumed clause \(D\) only
while its distinct witness clause \(C\) remains in the set.
One may also replace a clause by a condensation.
Rename different clauses' variables apart when representing the output.
\end{hbdefinition}

\begin{hbtheorem}{thm:p12-small-clause-simplify}{proved=annex:p12.small.clause.simplify}{The exact universally quantified clause guarantees}
Separating clause variable names, deleting tautologies, deleting a
subsumed clause with a retained witness, and replacing a clause by
a condensation preserve the models of the universally read clause set.
The deletion/condensation process terminates because clause count
plus literal-occurrence count strictly decreases.
Subsumption uses multiplicities: \(\{R(x,x),R(y,y)\}\)
does not subsume \(\{R(a,a)\}\), although it has the
equivalent condensation \(\{R(x,x)\}\).
Mutually subsuming clauses must not all be deleted simultaneously.
\end{hbtheorem}

\begin{hbtheorem}{thm:p12-small-pipeline}{proved=annex:p12.small.pipeline}{The complete standard transformation}
Apply existential closure when required, then simplification, NNF,
miniscoping, indexed rectification, outermost local Skolemization,
prefix formation and matrix distribution, and clause simplification.
The procedure terminates with a finite clause set \(S\) in a finite
signature extension. For each original structure \(\mathfrak A\):
\begin{enumerate}
\item Assuming Choice, if some valuation satisfies the input,
      \(\mathfrak A\) has an expansion satisfying \(S\).
\item Without Choice, if an expansion satisfies \(S\), some
      valuation satisfies the input in its reduct.
\end{enumerate}
For sentence inputs, replace \emph{some valuation satisfies} by
ordinary sentence satisfaction. In particular the procedure
preserves satisfiability, with the precise stronger model relation shown.
\end{hbtheorem}

\section{The complete running example and failure cases}

\begin{hbtheorem}{thm:p12-small-example}{proved=annex:p12.small.example}{Miniscoping, fresh variables and the resulting two clauses}
For the source's sentence
\[
 \forall x\exists y\forall z\,
       [R(x,x)\land(P(y)\lor R(x,y)\lor Q(z))],
\]
miniscoping yields
\[
 (\forall x R(x,x))\land
       [\exists y P(y)\lor\forall x\exists y R(x,y)\lor\forall z Q(z)].
\]
Indexed rectification uses five distinct binder names. Local
Skolemization uses a new constant \(a\) and a new unary \(f\).
The final clauses, with variables separated, are
\[
                     \{R(u,u)\},\qquad
                     \{P(a),R(v,f(v)),Q(w)\}.
\]
The proof supplies every intermediate formula. Without miniscoping,
one unary witness function suffices, but its argument also occurs
in the \(P\)-literal. Thus fewer arguments in some terms need not
mean fewer new function symbols.
\end{hbtheorem}

\begin{hbtheorem}{thm:p12-small-failures}{proved=annex:p12.small.failures}{Why the guards and the conjunctive structure matter}
Complete finite countermodels establish these failures:
dropping a miniscoping free-variable guard, distributing a universal
over disjunction or an existential over conjunction, raw renaming
that captures a free variable, treating a fresh-in-formula old
constant as a freely interpreted extension, deleting every member
of a subsumption cycle, and replacing the running example's two
clauses by the three pairwise disjunctions
\[
 \{R(u,u),P(a)\},\quad
 \{R(v,v),R(t,f(t))\},\quad
 \{R(w,w),Q(z)\}.
\]
The last replacement can turn an unsatisfiable problem, after adding
the opposite diagonal clause, into a satisfiable one.
\end{hbtheorem}

\begin{hbmethod}{def:p12-small-use}{What to supply and what to return}
Supply the finite formula tree, its signature, the intended open-variable
convention and an effective fresh-name supply. Return the clause set,
the new-symbol/dependency record and the sequence of transformations.
For validity checking, apply satisfiability-preserving clausification
to the negation of the sentence and use an appropriate refutation
procedure. The clause transformation does not decide arbitrary
first-order satisfiability by itself.
The optimized naming, witness and simplification methods preserve
their own source-specific conditions and cost claims.
\end{hbmethod}
